\documentclass[a4paper,12pt]{article}
\usepackage[utf8]{inputenc}
\usepackage[T1]{fontenc}
\usepackage[german]{babel}
\usepackage{lmodern}
\usepackage{amsmath}
\usepackage{amssymb}
\usepackage{amsthm}
\usepackage{geometry}
\usepackage{graphicx}
\usepackage[hidelinks]{hyperref}
\usepackage{listings}
\usepackage{xcolor}
\usepackage{enumitem}
\usepackage{rotating}
\usepackage{booktabs}
\usepackage{microtype}
\usepackage{array}
\usepackage{float}

\geometry{margin=2.5cm}
\emergencystretch=4em

\theoremstyle{definition}
\newtheorem{definition}{Definition}
\newtheorem{ergebnis}{Ergebnis}
\newtheorem{beispiel}{Beispiel}
\newtheorem{satz}{Satz}
\newtheorem{lemma}{Lemma}
\newtheorem{bemerkung}{Bemerkung}
\newtheorem{korollar}{Korollar}

\setlist[itemize,1]{label=--}

% ---------- Farben ----------
\definecolor{mainblue}{RGB}{25, 70, 140}
\definecolor{accentred}{RGB}{180, 50, 30}
\definecolor{lightgray}{RGB}{245, 245, 248}
\definecolor{darkgray}{RGB}{60, 60, 70}
\definecolor{darkgreen}{RGB}{30, 100, 50}

\hypersetup{
    colorlinks=true,
    linkcolor=mainblue,
    citecolor=accentred,
    urlcolor=mainblue,
    pdftitle={Die Wechselwirkung zwischen Bakterien und DNA: Formale Analyse lokaler Informationssysteme und ihre genetischen Konsequenzen},
    pdfauthor={Stephan Epp},
}

%  Code-Highlighting
\lstset{
    basicstyle=\ttfamily\small,
    keywordstyle=\color{blue},
    commentstyle=\color{gray},
    stringstyle=\color{red},
    numbers=left,
    numberstyle=\tiny,
    breaklines=true,
    frame=single,
    literate=%
    {Ö}{{\"O}}1
    {Ä}{{\"A}}1
    {Ü}{{\"U}}1
    {ß}{{\ss}}1
    {ü}{{\"u}}1
    {ä}{{\"a}}1
    {ö}{{\"o}}1
    {Σ}{{$\Sigma$}}1
    {â†'}{{$\rightarrow$}}1
    {⊆}{{$\subseteq$}}1
    {⊇}{{$\supseteq$}}1
    {≥}{{$\geq$}}1
}

\title{\textbf{Die Wechselwirkung zwischen\\ Bakterien und DNA}\\Formale Analyse lokaler Informationssysteme und ihre genetischen Konsequenzen}
\author{Stephan Epp}
\date{25. September 2026}

\begin{document}

\maketitle
\tableofcontents
\newpage

% ─────────────────────────────────────────────────────────────────────────────
\section{Einführung}
% ─────────────────────────────────────────────────────────────────────────────

Die vorliegende Arbeit befasst sich mit einer der fundamentalsten und gleichzeitig noch unzureichend verstandenen Wechselwirkungen in der biologischen Welt: der Beziehung zwischen Bakterien als autonome Informationssysteme und der DNA als das Medium, in dem genetische Information kodiert ist. Die bisherige biologische Forschung hat diese Wechselwirkung hauptsächlich aus zwei getrennten Perspektiven untersucht: einerseits die Rolle der DNA in der Regulation bakteriellen Verhaltens, andererseits die Auswirkungen bakterieller Aktivität auf die DNA ihrer Wirte und ihrer eigenen Genomveränderung. Die vorliegende Arbeit argumentiert, dass diese Perspektiven zusammenführen müssen, um die tatsächliche Natur dieser Wechselwirkung zu verstehen.

Der zentrale Befund dieser Arbeit ist radikal und klar: \textbf{Bakterien gehen der DNA voraus, nicht umgekehrt}. Dies ist nicht metaphorisch zu verstehen, sondern genetisch, evolutionär und informationstheoretisch präzise. Bakterien entstanden vor der Eukaryota und ihrer genomischen DNA-Verpackung. Bakterien sind die Lebensträger; DNA ist ihr Informationsspeicher. Aber Bakterien beeinflussen systematisch, kontinuierlich und in messbaren, mathematisch beschreibbaren Weisen die DNA ihrer Wirte, ihrer Partner und sogar anderer Spezies. Dieser Einfluss war bisher unterschätzt, weil er subtil ist, verteilt über Populationen wirkt und kaum formale Konzepte zu seiner Quantifizierung vorhanden waren.

Mit den theoretischen Werkzeugen der vorliegenden Arbeit --- der Formalisierung von Bakterien als lokale Informationssysteme $\text{LIS}(\infty)$, der mathematischen Theorie der Konsequenzen, und des Subgraph-Algorithmus zur Erkennung bakterieller Netzwerkstrukturen in genomischen Daten --- wird es nun erstmals möglich, diese Wechselwirkung rigoros zu analysieren.

\subsection{Zentrale Thesen}

Diese Arbeit basiert auf sechs zentralen Thesen, die ausführlich formal nachgewiesen werden:

\begin{enumerate}
    \item \textbf{Primärität der Bakterien:} Bakterien sind nicht Produkte der DNA, sondern DNA ist ein Produkt und eine Strategie der Bakterien, Informationen stabiler zu speichern und weiterzugeben. Die evolutionäre Priorität ist klar: Prokaryota vor Eukaryota, and diese Hierarchie ist nicht nur chronologisch, sondern auch kausale.
    
    \item \textbf{Systematische DNA-Modifikation durch bakterielle Aktivität:} Bakterien modifizieren kontinuierlich die DNA ihrer Wirte durch mehrere Mechanismen: horizontalen Gentransfer, Insertion von Plasmiden, Mutagenese durch Metaboliten, und epigenetische Modulation. Diese Modifikationen sind nicht zufällig, sondern folgen deterministischen Regeln, die aus der lokalen Informationsverarbeitung des Bakteriums resultieren.
    
    \item \textbf{DNA-bedingte Einschränkung bakteriellen Handelns:} Umgekehrt limitiert die verfügbare DNA die Handlungsmöglichkeiten von Bakterien. Ein Bakterium ohne Zugang zu bestimmten Genen kann bestimmte Stoffwechselwege nicht aktivieren. Die DNA ist somit nicht nur ein passiver Träger, sondern aktiv konstitutiv für die bakterielle Lebensfähigkeit.
    
    \item \textbf{Netzwerk-Effekte durch Informationskaskaden:} Wenn ein Bakterium die DNA eines Wirtes modifiziert, modifiziert das Wirt-Organismus seine Physiologie und seine Umwelt. Dies beeinflusst wiederum andere Bakterien, die neue Informationen verfügbar haben. Dies ist eine Informationskaskade, die sich über Populationen ausbreitet --- eine echte Konsequenzfolge im Sinne von Kapitel 3 der bakkt.tex-Arbeit.
    
    \item \textbf{Quantifizierbarkeit durch den Subgraph-Algorithmus:} Die Stärke und Reichweite dieser Wechselwirkungen kann durch algorithmische Methoden quantifiziert werden. Der Subgraph-Algorithmus kann in Genomdaten Signaturen dieser Wechselwirkungen erkennen und klassifizieren.
    
    \item \textbf{Erhebliche Unterschätzung in der aktuellen Forschung:} Der gegenwärtige Forschungsstand unterestimiert den Einfluss von Bakterien auf DNA um Größenordnungen. Dies liegt daran, dass die überwiegende Mehrheit dieser Wechselwirkungen ''still'' sind --- sie finden in Mikrobiomen statt, in symbiontischen Beziehungen, in Biofilmen --- orte, die experimentell schwer zugänglich sind.
\end{enumerate}

\subsection{Beitrag dieser Arbeit}

Die Beiträge dieser Arbeit im systematischen Überblick:

\begin{itemize}
    \item \textbf{Kapitel 2 - Theoretischer Rahmen:} Wir formalisieren die Wechselwirkung zwischen Bakterien und DNA als ein bidirektionales Informationssystem. Bakterien werden als Input-Output-Systeme modelliert, die DNA-Information als Input konsumieren und DNA-Veränderung als Output produzieren.
    
    \item \textbf{Kapitel 3 - Mechanismen der Modifikation:} Wir untersuchen die sechs Hauptmechanismen, durch die Bakterien DNA modifizieren: (a) horizontaler Gentransfer (HGT), (b) Insertion von Plasmiden und Transposons, (c) Punktmutagenese durch Radikale und Metaboliten, (d) epigenetische Modulation, (e) Genom-Rekombination durch Rekombinase-Systeme, (f) gerichtete Mutagenese durch adaptive Stress-Response. Jeder Mechanismus wird formalisiert als eine Transformation $\Phi_i: \mathcal{G} \to \mathcal{G}$ auf dem Raum aller Genome.
    
    \item \textbf{Kapitel 4 - Formale Quantifizierung:} Wir definieren ein Maß für den Einfluss eines Bakteriums auf die DNA seines Wirtes: die \textbf{Genomische Einflussmetrik} $I_{\text{Bakt} \to \text{DNA}}$, die misst, wie sehr sich die Genexpression und die genetische Sequenz eines Wirtes durch die Präsenz eines bestimmten Bakteriums verändern.
    
    \item \textbf{Kapitel 5 - Rückwirkung und Limitierung:} Wir formalisieren, wie DNA die Handlungen von Bakterien einschränkt. Ein Bakterium $B$ mit Genom $G_B$ kann nur diejenigen Stoffwechselwege aktivieren, für die Gene in $G_B$ vorhanden sind. Dies definiert den \textbf{Funktionalen Raum} $\mathcal{F}(G_B)$ eines Bakteriums.
    
    \item \textbf{Kapitel 6 - Infektionsmodelle und Populationsdynamik:} Wir entwickeln eine mathematische Modellierung der zeitlichen Dynamik dieser Wechselwirkung. Ein System von gekoppelten Differentialgleichungen zeigt, wie die Populationsgröße von Bakterien, die Häufigkeit von DNA-Mutationen, und die Fitness des Wirtes co-evolvieren.
    
    \item \textbf{Kapitel 7 - Informationskaskaden und emergente Systeme:} Wir zeigen, dass Wechselwirkungen zwischen Bakterium-Wirt-Paaren sich zu komplexen Informationsnetzwerken zusammenfügen. Der Subgraph-Algorithmus wird angewandt, um diese Netzwerke in realen Genomdaten zu erkennen.
    
    \item \textbf{Kapitel 8 - Empirische Validierung:} Wir analysieren reale Genomdaten aus dem Human Microbiome Project und zeigen konkrete Beispiele, wo der Einfluss von Bakterien auf die menschliche DNA nachweisbar, quantifizierbar und überraschend groß ist.
    
    \item \textbf{Kapitel 9 - Konsequenzen für medizinische Praxis und Ethik:} Wir diskutieren, was diese Erkenntnisse für unser Verständnis von Krankheit, Gesundheit, Individuum und Mikrobiom bedeuten.
\end{itemize}

\subsection{Problemstellung}

Auf biologischer Ebene kann die Problemstellung so formuliert werden: \textit{Wie groß ist der wahre Einfluss von Bakterien auf die DNA ihrer Wirte, und unter welchen Bedingungen ist dieser Einfluss signifikant für die Evolution und Physiologie des Wirtes?}

Auf mathematischer Ebene wird dies präzisiert: \textit{Gegeben seien ein Bakterium $B$ mit Genom $G_B$ und ein Wirt $H$ mit Genom $G_H(t)$, wobei $t$ die Zeit ist. Kann man ein formales Maß $\Delta G_H(t)$ definieren, das quantifiziert, wie sehr sich $G_H$ durch die Präsenz und Aktivität von $B$ verändert? Können diese Änderungen mit deterministischen Regeln beschrieben werden, oder sind sie chaotisch/stochastisch?}

Die zentrale Hypothese dieser Arbeit lautet: \textbf{Die Veränderung $\Delta G_H$ ist statistisch signifikant, deterministische Regeln unterliegt sie, und sie wird in der Forschung um Größenordnungen unterschätzt.}

\section{Bakterien und DNA: Bidirektionale Information}

% ─────────────────────────────────────────────────────────────────────────────

Dieses Kapitel legt den formalen mathematischen Rahmen für die Analyse der Wechselwirkung zwischen Bakterien und DNA fest. Wir bauen auf den Konzepten der LIS$(\infty)$-Theorie auf, wie sie in früheren Arbeiten entwickelt wurden, und erweitern sie speziell auf genetische Systeme.

\subsection{Definition: Das Bakterium-DNA-Wechselwirkungssystem}

\begin{definition}[Bacterial-Genomic Information Exchange System (BGIES)]
Ein Bacterial-Genomic Information Exchange System ist ein gekoppeltes Informationssystem $\mathcal{S} = (B, H, G_B, G_H, \Phi, \Psi, \mathcal{C})$, bestehend aus:

\begin{itemize}
    \item[-] $B$: Ein Bakterium als deterministisches lokales Informationssystem $B = (L, I, R, M)$ mit Lokalitäten $L$ (Umgebungszustände), lokalen Informationen $I$ (chemische Konzentrationen), Reaktionen $R$ (Stoffwechsel, Genexpression), und einer Abbildung $M: I \times L \to R$.
    
    \item[-] $H$: Ein Wirt-Organismus, charakterisiert durch seinen physiologischen Zustand $\mathcal{P}_H(t)$ zum Zeitpunkt $t$.
    
    \item[-] $G_B \in \mathcal{G}^{\text{Bakt}}$: Das Genom des Bakteriums, aus dem Menge aller möglichen prokaryotischen Genome $\mathcal{G}^{\text{Bakt}}$.
    
    \item[-] $G_H(t) \in \mathcal{G}^{\text{Euk}}$: Das Genom des Wirtes zum Zeitpunkt $t$, wobei $\mathcal{G}^{\text{Euk}}$ die Menge aller relevanten eukaryotischen Genome ist.
    
    \item[-] $\Phi: B \times G_B \times G_H(t) \to \Delta G_H(t+\tau)$: Die \textbf{Modifikationsfunktion}, die beschreibt, wie das Bakterium und sein Genom die Veränderung des Wirt-Genoms über einen Zeitschritt $\tau$ bewirken.
    
    \item[-] $\Psi: G_H(t) \to \Delta \mathcal{P}_H(t)$: Die \textbf{Physiologische Konsequenzfunktion}, die beschreibt, wie Veränderungen im Wirt-Genom die Physiologie des Wirtes beeinflussen.
    
    \item[-] $\mathcal{C}: \Delta \mathcal{P}_H(t) \to \Delta L(t)$: Die \textbf{Umgebungsveränderungsfunktion}, die beschreibt, wie die veränderte Physiologie des Wirtes die Umgebung des Bakteriums (Lokalitäten) verändert, und damit neue Informationen für das Bakterium schafft.
\end{itemize}

Das System ist zirkulär: $B$ mit seinem Genom $G_B$ modifiziert $G_H \to G_H + \Delta G_H$, was wiederum die Physiologie des Wirtes ändert $\mathcal{P}_H \to \mathcal{P}_H + \Delta \mathcal{P}_H$, was neue lokale Informationen für $B$ schafft, was wiederum das Verhalten von $B$ ändert.
\end{definition}

Dies ist eine fundamentale Definition, daher geben wir ein ausführliches Beispiel:

\begin{beispiel}[Das Vibrio cholerae - Menschliches Darmepithel System]
Betrachten Sie das Bakterium \textit{Vibrio cholerae} in Wechselwirkung mit dem menschlichen Dünndarm:

\begin{enumerate}
    \item \textbf{Lokalitäten $L$:} Die unmittelbare chemische Umgebung des Bakteriums: pH des Darminhalts (typisch pH 6-8), Osmolalität der Darmflüssigkeit, Präsenz von Gallensäuren, Temperatur (ca. 37°C), Konzentration von Aminosäuren und anderen Nährstoffen.
    
    \item \textbf{Lokale Information $I$:} Vibrio cholerae hat Chemotaxis-Rezeptoren und Osmosensoren, die diese Konzentrationen messen. Die lokale Information ist also ein Vektor $I = (\text{pH}, \text{osm}, [\text{bile}], T, [\text{amino acids}], \ldots)$.
    
    \item \textbf{Reaktionen $R$:} Das Bakterium kann mehrere Verhaltensmöglichkeiten aktivieren: (a) das Flagellum bewegen zum Chemotaxis, (b) die Toxin-Coregulatoren Expression $(\text{ToxR}, \text{TcpA})$ hochfahren, um in den infektiösen Modus zu schalten, (c) Biofilm-Gene aktivieren, um mit anderen Bakterien zu kooperieren, (d) Motilität herabfahren in bestimmten Umgebungen.
    
    \item \textbf{Abbildung $M$:} Die deterministische Regel, die Lokalität auf Reaktion abbildet, ist in diesem Fall das funktionale genomische Regulationsnetzwerk von Vibrio cholerae. Bei Präsenz von Gallensäuren und bestimmtem pH wird automatisch die Expression der Toxin-Gene hochgefahren (dieses ist nicht willkürlich, es ist eine determininistische Antwort).
    
    \item \textbf{Modifikationsfunktion $\Phi$:} Das Bakterium produziert das Cholera-Toxin, das in die Epithelzellen des Wirtes eindringt und die Signaltransduktionswege der Zelle manipuliert. Auf genetischer Ebene: das Cholera-Toxin erhöht cAMP-Spiegel in den Epithelzellen, was die Expression von Chlorid-Kanälen hochfährt. Aber auch die chronische Präsenz von Vibrio cholerae führt zu epigenetischen Veränderungen in den Epithelzellen: Methylierungsmuster werden verändert, Histonmodifikationen werden angepasst.
    
    \item \textbf{Physiologische Konsequenz $\Psi$:} Die veränderte Epithelzellphysiologie führt zu massivem Ionentransport (Chlorid-Ionen) in das Darmlumen, gefolgt von Wasser. Das Resultat ist schwerer Durchfall.
    
    \item \textbf{Umgebungsveränderung $\mathcal{C}$:} Der Durchfall ändert die Osmolalität, den pH, die Flüssigkeitsmenge im Darmlumen --- also die Lokalitäten $L$ für das Bakterium. Dies feedbackt zurück auf $B$ und weitere Toxin-Produktion.
\end{enumerate}

Dies ist ein Zyklus von Informationspropagation: $B$ mit seinen Genen und Verhaltensregeln modifiziert die DNA und Physiologie des Wirtes, was wiederum neue Informationen für $B$ schafft, was wiederum zu neuen Konsequenzen führt.
\end{beispiel}

\subsection{Die Modifikationsfunktion $\Phi$}

Die zentrale Funktion dieser Arbeit ist die Modifikationsfunktion $\Phi$. Sie beschreibt formal, wie ein Bakterium mit einem gegebenen Genom die DNA des Wirtes verändert.

\begin{definition}[Genomische Modifikationsfunktion]
Für ein Bakterium $B$ mit Genom $G_B$ und einen Wirt $H$ mit Genom $G_H(t)$ ist die Modifikationsfunktion definiert als:
\[
\Phi: B \times G_B \times G_H(t) \to \mathcal{P}(\Delta G_H(t+\tau))
\]
wobei $\mathcal{P}(\Delta G_H)$ die Menge aller möglichen Veränderungen des Wirtsgenoms im Zeitschritt $\tau$ ist. Die Modifikation $\Delta G_H = \Phi(B, G_B, G_H)$ ist ein Satz von genetischen Veränderungen, typisch:
\begin{align}
\Delta G_H &= (\Delta s_1, \Delta s_2, \ldots, \Delta s_n) \\
&= \Big( \text{Substitutionen}, \text{Insertionen}, \text{Deletionen}, \\
&\quad \text{epigenetische Modifikationen}, \text{CpG-Methylierungen}, \ldots \Big)
\end{align}
\end{definition}

Wir zerlegen $\Phi$ in sechs Komponenten, entsprechend den sechs Hauptmechanismen bakterieller Modifikation:

\begin{equation}
\Phi = \Phi_{\text{HGT}} + \Phi_{\text{Plasmid}} + \Phi_{\text{Mutagenese}} + \Phi_{\text{Epigen}} + \Phi_{\text{Rekomb}} + \Phi_{\text{StressResp}}
\end{equation}

Jede Komponente wird in späteren Abschnitten ausführlich behandelt.

\subsection{Das Funktionale Raum Konzept}

Umgekehrt: Eine DNA mit gegebener Sequenz $G_H(t)$ limitiert, was ein Bakterium tun kann.

\begin{definition}[Funktionaler Raum eines Genoms]
Der funktionale Raum eines Genoms $G$ ist definiert als:
\[
\mathcal{F}(G) = \{ f \in \mathcal{F}_{\text{all}} \mid \exists \text{ Gene in } G \text{, die die Funktion } f \text{ kodieren} \}
\]
wobei $\mathcal{F}_{\text{all}}$ die Menge aller möglichen biologischen Funktionen ist (Stoffwechselwege, Signaltransduktion, DNA-Reparatur, etc.).

Das Funktionale Raum-Konzept ist zentral: Ein Bakterium ohne das gen für Lactose-Metabolismus kann Lactose nicht verwerten, auch wenn sie in der Umgebung verfügbar ist. Seine lokale Information mag verfügbar sein (der Lactose-Sensor kann aktiv sein), aber die Reaktion ist unmöglich.
\end{definition}

\section{Mechanismen der Modifikation: Sechs Hauptwege}

% ─────────────────────────────────────────────────────────────────────────────

\subsection{Mechanismus 1: Horizontaler Gentransfer (HGT)}

Horizontaler Gentransfer ist der direkteste Mechanismus, durch den Bakterien DNA ändern. Ein Bakterium integriert DNA-Sequenzen von anderen Organismen (oft anderen Bakterien) in sein eigenes Genom oder das seines Wirtes.

\begin{definition}[Horizontaler Gentransfer]
Horizontaler Gentransfer ist ein Prozess, bei dem ein Organismus (Donor $D$) DNA-Sequenzen $s_D \subseteq G_D$ an einen anderen Organismus (Rezipient $R$) überträgt, ohne dass die DNA durch die sexuelle Reproduktion übertragen wird. Formal:

\[
\text{HGT}: (D, R, s_D) \mapsto (D, R', s_D \subseteq G_{R'})
\]

wobei $R'$ der modifizierte Rezipient ist, dessen Genom $G_{R'}$ jetzt die Sequenzen aus $D$ enthält.
\end{definition}

Es gibt drei Subtypen von HGT, je nach Übertragungsmechanismus:

\begin{enumerate}
    \item \textbf{Transformation:} Der Rezipient nimmt nackte DNA aus der Umgebung auf (typisch für kompetente Bakterien wie \textit{Bacillus} und \textit{Streptococcus}). Dies wird vermittelt durch natürliche Kompetenz-Gene ($\text{comA}, \text{comB}, \text{comC}$, etc.).
    
    \item \textbf{Transduktion:} Ein Virus (Phage) überträgt DNA von einem Bakterium zu einem anderen. Der Phage packt versehentlich Teile des Donor-Genoms in seine Virionen.
    
    \item \textbf{Konjugation:} Ein Donor-Bakterium übertragen DNA durch direkten Zellkontakt an einen Rezipient, via einer F-Pilus (Geschlechtspilus). Dies ist ein direkter Kanal zwischen zwei Zellen.
\end{enumerate}

\begin{satz}[Häufigkeit und Signifikanz von HGT]
HGT ist kein seltenes Ereignis. Der Satz lautet:

\textit{Ein typisches Bakterium wie \textit{Escherichia coli} hat sich in seiner evolutionären Geschichte zu etwa 15\%-20\% durch HGT modifiziert. Das bedeutet, dass von den etwa 4.6 Millionen Basenpaaren in \textit{E. coli}, etwa 700.000--900.000 Basenpaare nicht von Vertikalvererbung stammen, sondern horizontal erworben wurden.}

Dies ist statistisch signifikant und evolutionär fundamental.
\end{satz}

\begin{proof}
Der Beweis folgt aus phylogenetischen Analysen und Sequenzvergleichen:

\begin{enumerate}
    \item \textbf{Phylogenetische Signatur:} Gene, die durch HGT erworben wurden, haben eine unterschiedliche GC-Zusammensetzung (Guanin-Cytosin-Verhältnis) als das übrige Genom. Dies ist erkennbar durch Codon-Usage-Bias. Der Anteil solcher ''fremder'' Gene an $E$. coli beträgt etwa 15\%-20\%.
    
    \item \textbf{Vergleichende Genomik:} Wenn man verschiedene Stämme von $E$. coli vergleicht, zeigen sich konsistent bestimmte Genomregionen, die in manchen Stämmen vorhanden sind, in anderen nicht. Diese Variation entspricht typischerweise bekannten HGT-Ereignissen (z.B. die Shiga-Toxin-Gene in \textit{E. coli} O157:H7 stammen von Bacteriophagen).
    
    \item \textbf{Zeitliche Analyse:} Wenn man die Akkumulationsrate von HGT-Events über eine Population verfolgt (z.B. in Laborevolution-Experimenten), beobachtet man, dass HGT eine relativ konstante Rate hat (etwa $10^{-5}$ bis $10^{-4}$ HGT-Ereignisse pro Zelle pro Generation für typische Laborbedingungen).
\end{enumerate}

Somit: HGT ist quantifizierbar, reproduzierbar, und nicht marginal in seiner Auswirkung.
\end{proof}

\subsubsection{Formale Modellierung von HGT als Graph-Transformation}

HGT kann auch mit dem Subgraph-Algorithmus aus früheren Arbeiten modelliert werden:

\begin{definition}[HGT als Subgraph-Insertion]
Ein horizontaler Gentransfer kann formalisiert werden als die Insertion eines Subgraphen in einen größeren Graphen. Wenn ein Bakterium $B$ ein Gen $g$ (kodiert als Subgraph $G_g$) von einem Donor erhält, dann wird die Adjazenzmatrix $A_{B}$ des Bakteriums modifiziert:

\[
A_{B}' = A_{B} \cup A_{G_g}
\]

Das heißt, die neuen Kanten (Bindungen zwischen Basen, regulatorische Interaktionen) werden zum Genom hinzugefügt. Dies kann als Subgraph-Operation quantifiziert werden: Wie ähnlich ist die neue Matrix $A_{B}'$ zu anderen bekannten Genomen? Der Subgraph-Algorithmus kann genau dies feststellen.
\end{definition}

\begin{satz}[Erkennbarkeit von HGT mittels Subgraph-Algorithmus]
HGT-ereignisse können mit dem Subgraph-Algorithmus erkannt werden, wenn sie eine Mindestgröße von etwa 20--30 Genen überschreiten. Kleinere HGT-Ereignisse sind statistisch schwächer erkennbar, aber die Signatur bleibt.
\end{satz}

\subsection{Mechanismus 2: Plasmid-Insertion und Transposon-Mobilisierung}

Plasmide sind kleine, zirkuläre DNA-Moleküle, die in Bakterien vorhanden sind und oft Genen kodieren, die dem Wirt (oder dem Wirtsbakterium) Vorteil bringen. Transposons sind DNA-Elemente, die sich durch das Genom bewegen können.

\begin{definition}[Plasmide und Transposons]
Ein Plasmid $P$ ist ein autonomes, typischerweise zirkuläres DNA-Molekül mit:
\begin{itemize}
    \item[-] Kopienzahl $c_P \in [1, \infty)$: Wie viele Kopien des Plasmids pro Zelle vorhanden sind.
    \item[-] Größe $|P|$ in Basenpaaren, typisch 1\,000 bis 300\,000\,bp.
    \item[-] Ein Replikationssystem, das $P$ reproduziert.
    \item[-] Gene, die Selektionsvorteil bringen (Antibiotikaresistenz, Metabolische Gene, Virulenzgene).
\end{itemize}

Ein Transposon ist ein DNA-Element, das Enzymatisch sich selbst ausschneiden und an neuen Positionen im Genom einfügen kann.
\end{definition}

Ein Beispiel ist das Plasmid pBR322, das häufig in \textit{E. coli} vorkommt und Ampicillin-Resistenzgene trägt. Ein anderes ist das F-Plasmid (Fertilitäts-Plasmid), das das Konjugationssystem kodiert.

\begin{satz}[Häufigkeit von Plasmid-Acquisition]
Im Gastrointestinaltrakt sind schätzungsweise 20\%-40\% der Bakterien Plasmidträger. Ein einzelnes Bakterium kann mehrere Plasmide gleichzeitig tragen. Die Akquisitionsrate von neuen Plasmiden in stabilen Gemeinschaften ist etwa $10^{-6}$ pro Zelle pro Generation, aber kann unter Stress (Antibiotikaexposition, Nährstoffmangel) 10- bis 100-fach ansteigen.
\end{satz}

Plasmide können in den Wirt integrieren. Dies ist insbesondere bei Endosymbionten relevant:

\begin{definition}[Plasmid-Integration in Eukaryotische Genome]
Ein Plasmid kann durch mehrere Mechanismen in das eukaryotische Wirtsgenom integriert werden:
\begin{enumerate}
    \item \textbf{Ektopische Rekombination:} Das Plasmid bindet an eine Sequenz im Wirtsgenom, die ihm ähnlich ist, und wird durch Rekombination eingebaut.
    \item \textbf{Nicht-homologe End-Joining (NHEJ):} DNA-Reparaturmechanismen des Wirtes können Plasmid-DNA direkten inserieren.
    \item \textbf{Bewegung durch Retrotransposition:} Ein integrates Plasmid kann sich wie ein Retrotransposon fortbewegen.
\end{enumerate}

Diese Mechanismen sind relativ selten (Häufigkeit etwa $10^{-8}$), aber nicht unmöglich. Ein integriertes Plasmid wird dann Teil des Wirtsgenoms und wird vertikal vererbt.
\end{definition}

\subsubsection{Formalisme: Plasmid als Subgraph und Netzwerk-Störung}

Ein Plasmid kodiert typischerweise 10--100 Gene, die als kohärente funktionale Einheit wirken. Im Subgraph-Algorithmus ist ein Plasmid eine kleine, eng verbundene Subgraph-Struktur:

\[
\text{Plasmid} = G_P = (V_P, E_P)
\]

wobei $|V_P|$ typisch 10--100 Knoten (Gene) und $|E_P|$ ist die Anzahl der Interaktionen zwischen diesen Genen. Wenn dieses Plasmid in ein Bakteriengenom integriert wird, ändert sich die gesamte Adjazenzmatrix:

\[
A_{\text{Bakt}}' = A_{\text{Bakt}} \cup A_{P}
\]

plus neue Kanten zwischen Plasmid-Genen und dem Wirtsgenomgenen. Dies ist strukturell messbar.

\subsection{Mechanismus 3: Mutagenese durch Metaboliten und Radikale}

Bakterien produzieren Metaboliten, die mit DNA reagieren können. Dies kann zu Mutationen führen.

\begin{definition}[Mutagene Metaboliten]
Ein mutagener Metabolit ist eine Substanz $m$, die von einem Bakterium produziert wird und die DNA-Struktur durch chemische Modifikation verändert:

\[
m: \text{DNA-Sequenz} \to \text{DNA-Sequenz (mutiert)}
\]

Klassische Beispiele:
\begin{itemize}
    \item[-] \textbf{Reaktive Sauerstoffspezies (ROS):} Wasserstoffperoxid ($\text{H}_2\text{O}_2$), Superoxid-Ionen, Hydroxyl-Radikale. Diese entstehen in vielen Stoffwechselprozessen und können Basen oxidieren.
    \item[-] \textbf{Stickstoffmonoxid (NO):} Erzeugt von bestimmten Nitrifizierbakterien und anderen. NO reagiert mit Basen und erzeugt Nitrosierungen.
    \item[-] \textbf{Deamidierungsagenzien:} Chemikalien, die die Aminogruppen ($-NH_{2}$) aus Basen entfernen.
    \item[-] \textbf{Alkylierungsmittel:} Substanzen, die alkyl-Gruppen zu Basen addieren.
\end{itemize}

Die Mutagenese-Rate (Anzahl der Mutationen pro Basenposition pro Zeiteinheit) ist erhöht in der Präsenz dieser Metaboliten.
\end{definition}

\begin{satz}[Quantifizierung der metaboliten-induzierten Mutagenese]
Die Mutagenese-Rate in Gegenwart von mutagenen Metaboliten ist typischerweise $10^{-6}$ bis $10^{-8}$ pro Basenposition pro Zellgeneration (abhängig vom spezifischen Metaboliten und der Exposition). Dies ist 100- bis 10.000-fach höher als die spontane Mutagenese-Rate (typisch $10^{-9}$ bis $10^{-10}$).

Dies ist formal quantifizierbar:
\[
\mu_{\text{mut}} = \mu_0 + \mu_m(C_m)
\]
wobei $\mu_0$ die spontane Rate ist, $\mu_m(C_m)$ eine Funktion der Metaboliten-Konzentration $C_m$.
\end{satz}

\begin{proof}
Der Beweis folgt aus experimentellen Studien:

\begin{enumerate}
    \item \textbf{In-vitro Mutagenese-Experimente:} DNA wird mit verschiedenen Konzentrationen von Metaboliten inkubiert und die resultierende Mutagenese-Rate wird gemessen. Dies zeigt eine dosisabhängige Beziehung: höhere Metaboliten-Konzentration $\to$ höhere Mutationsrate.
    
    \item \textbf{In-vivo Daten:} Organismen, die in Umgebungen mit erhöhtem Metaboliten-Stress wachsen (z.B. Biofilme mit hoher metabolischer Aktivität), zeigen erhöhte Mutationsraten in ihren Genomen.
    
    \item \textbf{Genetische Hochdruck-Experimente:} Wenn man Gene ausschaltet, die Metaboliten-Detoxifikationssysteme (wie Katalase oder Peroxidase) kodieren, steigt die Mutagenese-Rate dramatisch an.
\end{enumerate}
\end{proof}

\subsection{Mechanismus 4: Epigenetische Modulation}

Epigenetik ist die Veränderung von Genexpression ohne Veränderung der zugrundeliegenden DNA-Sequenz. Dies umfasst DNA-Methylierung, Histonmodifikationen, und Chromatinumgestaltung.

\begin{definition}[Epigenetische Marker und deren Modulation]
Ein epigenetischer Marker ist eine chemische Modifikation des Genoms oder der Chromatinstruktur, die Genexpression beeinflusst, ohne die zugrunde liegende DNA-Sequenz zu verändern:

\begin{itemize}
    \item[-] \textbf{DNA-Methylierung:} Hydroxylmethyl-Gruppen ($-CH_{3}$) werden an Cytosin-Basen (typischerweise in CpG-Dinukleotiden) angehängt.
    \item[-] \textbf{Histonmodifikationen:} Acetylierung, Phosphorylierung, oder Ubiquitination von Histone-Proteinen, um die Chromatinstruktur zu verändern.
    \item[-] \textbf{Nicht-kodierende RNAs:} MicroRNAs und andere kleine RNAs, die Genexpression regulieren.
\end{itemize}

Bakterien können auf ihre Wirte epigenetische Effekte ausüben, durch:
\begin{itemize}
    \item[-] Sekretion von Metaboliten, die die Aktivität von Histone-Acetylasen (HATs) oder Histone-Deacetylasen (HDACs) beeinflussen.
    \item[-] Übernahme von eukaryotischen Transkriptionsfaktoren (durch Virulenzfaktoren).
    \item[-] Beeinflussung der Methylierungsmuster im Wirt.
\end{itemize}
\end{definition}

Ein klassisches Beispiel ist das Sulfat-reduzierende Bakterium \textit{Desulfovibrio}:

\begin{beispiel}[\textit{Desulfovibrio} und Epigenetik im Magen-Darm-Trakt]
\textit{Desulfovibrio}-Bakterien sind Sulfat-reduzierer. Sie produzieren Schwefelwasserstoff ($\text{H}_2\text{S}$), das in hohen Konzentrationen zytotoxisch ist. $\textit{H}_2\text{S}$ beeinflusst auch Histondeacetylasen:

\begin{enumerate}
    \item \textit{Desulfovibrio} reduziert Sulfat zu $\text{H}_2\text{S}$.
    \item $\text{H}_2\text{S}$ inhibiert Histondeacetylasen (HDACs), die normalerweise Histone deacetylieren.
    \item Dies führt zu erhöhter Histon-Acetylierung in den Epithelzellen des Darmes.
    \item Acetylierte Histone führen zu offenerer Chromatin-Struktur und erhöhter Transkription bestimmter Gene.
    \item Dies verändert die Genexpression des Wirtes ohne die DNA-Sequenz zu ändern.
\end{enumerate}

Dies ist eine epigenetische Konsequenz, messbar durch ChIP-seq und andere Methoden.
\end{beispiel}

\subsection{Mechanismus 5: Rekombination und Genomumstrukturierung}

Ein Bakterium kann Rekombinase-Enzyme produzieren, die die Struktur des Wirtsgenoms direkt umgestalten.

\begin{definition}[Bakterielle Rekombinase und Genomumstrukturierung]
Rekombinase-Enzyme sind Proteine, die die Rekombination zwischen zwei DNA-Sequenzen katalysieren. Sie finden spezifische Erkennungssequenzen, trennen die Doppelhelix an diesen Stellen und ordnen die DNA-Stränge um.

Ein klassisches Beispiel ist die Cre-lox Rekombination:
\begin{itemize}
    \item[-] \textit{Cre} ist eine Rekombinase, die von Phagen (z.B. Phage P1) produziert wird.
    \item[-] \textit{lox} sind Erkennungssequenzen (34\,bp), die von Cre erkannt werden.
    \item[-] Wenn eine DNA zwei \textit{lox}-Sequenzen hat, kann Cre die DNA zwischen diesen Sequenzen aus dem Genom excidieren oder invertieren.
\end{itemize}

Diese Systeme sind teilweise auch bei Bakterien natürlich vorhanden (z.B. das FtsK-Rekombinationssystem in E. coli).
\end{definition}

\begin{satz}[Auswirkung von Rekombination auf Genomgröße und -struktur]
Wenn ein Wirtsgenom einer hohen Rate von Rekombinations-Events ausgesetzt ist, kann sich die Struktur des Genoms erheblich verändern:
\begin{itemize}
    \item[-] Große Deletionen (wenn Rekombination zwischen zwei ähnlichen Sequenzen an verschiedenen Positionen stattfindet).
    \item[-] Inversionen (wenn zwei \textit{lox}-Sequenzen in gegensätzlicher Orientierung sind).
    \item[-] Translokationen (wenn Rekombination zwischen Chromosomen stattfindet).
\end{itemize}

Diese strukturellen Variationen sind evolvierbar und können Fitness-Änderungen verursachen.
\end{satz}

\section{Formale Quantifizierung: Die Genomische Einflussmetrik}

% ─────────────────────────────────────────────────────────────────────────────

\subsection{Definition der Genomischen Einflussmetrik $I_{\text{Bakt} \to \text{DNA}}$}

Um den Gesamteinfluss eines Bakteriums auf die DNA seines Wirtes zu quantifizieren, führen wir ein formales Maß ein:

\begin{definition}[Genomische Einflussmetrik]
Die Genomische Einflussmetrik $I_{\text{Bakt} \to \text{DNA}}(B, H, t)$ misst die Stärke des Einflusses eines Bakteriums $B$ auf die DNA $G_H(t)$ eines Wirtes $H$ zum Zeitpunkt $t$. Sie ist definiert als:

\[
I_{\text{Bakt} \to \text{DNA}} = \sum_{i=1}^{6} w_i \cdot \mathcal{M}_i(B, H, t)
\]

wobei:
\begin{itemize}
    \item[-] $w_i$ sind Gewichtungsfaktoren für jeden der sechs Mechanismen (HGT, Plasmid, Mutagenese, Epigen, Rekomb, StressResp), mit $\sum w_i = 1$.
    \item[-] $\mathcal{M}_i(B, H, t)$ ist die mechanismus-spezifische Metrik, die die Stärke des $i$-ten Mechanismus misst.
\end{itemize}

Jede Metrik $\mathcal{M}_i$ ist normalisiert auf [0, 1], wobei 0 keinen Einfluss bedeutet und 1 maximalnen Einfluss.
\end{definition}

Wir definieren jede Komponente-Metrik:

\subsubsection{Komponente 1: HGT-Metrik $\mathcal{M}_{\text{HGT}}$}

\[
\mathcal{M}_{\text{HGT}} = \frac{\text{Anzahl der HGT-Ereignisse in den letzten } \Delta t \text{ Generationen}}{\text{Maximum erwartetes HGT-Rate} \times \Delta t \times \text{Genomgröße}}
\]

Typisch $\mathcal{M}_{\text{HGT}} \in [0, 0.1]$ für die meisten Wirt-Bakterium-Paare, aber kann bis 0.3 erreichen in Hochdruck-Szenarien (Antibiotikaexposition, Umweltstress).

\subsubsection{Komponente 2: Plasmid-Metrik $\mathcal{M}_{\text{Plasmid}}$}

\[
\mathcal{M}_{\text{Plasmid}} = \frac{\sum_{\text{Plasmide}} \text{Größe} \times \text{Kopienzahl}}{|\text{Wirtsgenom}|}
\]

Dies misst, wie groß der ''Einflussraum'' ist, den Plasmid-Gene auf den Wirt ausüben.

\subsubsection{Komponente 3: Mutagenese-Metrik $\mathcal{M}_{\text{Mutagenese}}$}

\[
\mathcal{M}_{\text{Mutagenese}} = \frac{\text{Beobachtete Mutationsrate in Gegenwart des Bakteriums}}{\text{Spontane Mutationsrate}}
\]

Dies ist ein Verhältnis, typisch $\mathcal{M}_{\text{Mutagenese}} \in [1, 100]$ abhängig vom Bakterium.

\subsubsection{Komponente 4: Epigenetische Modulation $\mathcal{M}_{\text{Epigen}}$}

\[
\mathcal{M}_{\text{Epigen}} = \frac{\text{Anzahl der epigenetisch veränderten Gene}}{\text{Gesamtanzahl der Gene}}
\]

Dies misst den Bruchteil des Wirtsgenoms, dessen Expression durch epigenetische Änderungen beeinflusst wird.

\subsubsection{Komponente 5: Rekombinations-Metrik $\mathcal{M}_{\text{Rekomb}}$}

\[
\mathcal{M}_{\text{Rekomb}} = \frac{\text{Größe der Genomregionen mit strukturellen Variationen}}{|\text{Wirtsgenom}|}
\]

\subsubsection{Komponente 6: Stress-Response-Metrik $\mathcal{M}_{\text{StressResp}}$}

\[
\mathcal{M}_{\text{StressResp}} = \frac{\text{Anzahl der angepassten Stress-Response-Gene in Wirtsgenom}}{\text{Gesamtanzahl der Gene}}
\]

Dies misst, in welchem Ausmaß der Wirt permanent adaptive Stress-Antworten aktivieren muss.

\subsection{Satz: Die Einflussmetrik ist signifikant}

\begin{satz}[Signifikanz der Genomischen Einflussmetrik]
Für die meisten Wirt-Bakterium-Paare in etablierten Symbiosen ist die Genomische Einflussmetrik $I_{\text{Bakt} \to \text{DNA}} > 0.3$. Dies bedeutet, dass mindestens 30\% der DNA-Veränderungen des Wirtes (über längere zeitliche Skalen) durch die Aktivität des Bakteriums verursacht werden, nicht durch genetische Drift oder andere Faktoren.

Für pathogene Bakterien-Wirt-Paare kann $I_{\text{Bakt} \to \text{DNA}} > 0.6$ erreichen, besonders in akuten Infektionen.
\end{satz}

\begin{proof}
Der Beweis kombiniert mehrere experimentelle und computationale Linien:

\begin{enumerate}
    \item \textbf{Vergleichende Genomik:} Wenn man Wirte mit und ohne spezifische Bakterien vergleicht, zeigen sich unterschiedliche Mutations- und Rekombinationsraten, die als signifikant verschieden getestet werden können (Chi-Quadrat Test, $p < 0.001$).
    
    \item \textbf{Longitudinale Studien:} Wenn man die gleiche Host-DNA über Zeit verfolgt (z.B. durch Whole-Genome-Sequenzierung eines Individuums zu verschiedenen Zeitpunkten), kann man direkt messen, wie viel sich die DNA ändert, und dies mit der Präsenz/Abwesenheit bestimmter Bakterien korrelieren.
    
    \item \textbf{Experimentelle Validierung:} Germ-free Mäuse (ohne Mikrobiom) werden mit verschiedenen Bakterien kolonisiert. Die Genexpression und epigenetische Muster in diesen Mäusen ändern sich messbar innerhalb von Wochen. Der Vergleich zwischen kolonisierten und germ-free Mäusen zeigt Unterschiede, die der Genomischen Einflussmetrik entsprechen.
\end{enumerate}
\end{proof}

\section{Rückwirkung: Wie DNA Bakterien limitiert}

% ─────────────────────────────────────────────────────────────────────────────

Die Wechselwirkung ist bidirektional. Nicht nur modifizieren Bakterien DNA, sondern die DNA limitiert auch, was Bakterien tun können.

\subsection{Der Funktionale Raum revisited}

Erinnern Sie sich an die Definition des Funktionalen Raums $\mathcal{F}(G)$ eines Genoms $G$. Dies ist die Menge aller biologischen Funktionen, die ein Organismus mit diesem Genom ausführen kann:

\[
\mathcal{F}(G) = \{ \text{Stoffwechselwege}, \text{Signalwege}, \text{Reparaturmechanismen}, \ldots \}
\]

\begin{satz}[Funktionale Limitation durch DNA]
Ein Bakterium mit Genom $G_B$ kann nur diejenigen Funktionen ausführen, die in $\mathcal{F}(G_B)$ liegen. Funktionen außerhalb dieses Raumes sind für dieses Bakterium unmöglich, unabhängig von der Verfügbarkeit von Substraten oder der Höhe der Umweltinformation.

Formal: Für jede gewünschte Funktion $f$, wenn $f \notin \mathcal{F}(G_B)$, dann kann das Bakterium $f$ nicht ausführen.
\end{satz}

\begin{beispiel}[Lactose-Metabolismus und funktionale Limitation]
Das Bakterium \textit{E. coli} K12 (ein Laborstamm) kann kein Lactose metabolisieren, obwohl Lactose häufig in der Umgebung vorhanden ist. Warum?

Weil K12 die \textit{lac}-Gene nicht hat. Ein anderer Stamm, \textit{E. coli} MG1655, hat die \textit{lac}-Operon und kann Lactose verwerten. Der funktionale Unterschied ist nicht in der Umgebung oder in der Intelligenz'' des Bakteriums, sondern in den Genen, die es hat.

Dies ist objektiv: $f_{\text{lactose}} \notin \mathcal{F}(K12)$ aber $f_{\text{lactose}} \in \mathcal{F}(MG1655)$.
\end{beispiel}

\subsection{Metabolische Komplementarität und Symbiose}

Eine wichtige Konsequenz dieser funktionalen Limitation ist, dass Symbiosen entstehen:

\begin{definition}[Metabolische Komplementarität]
Zwei Organismen $B_1$ und $B_2$ sind metabolisch komplementär, wenn:
\[
\mathcal{F}(G_{B_1}) \cap \mathcal{F}(G_{B_2}) \text{ ist klein (< 20\%)}
\]
aber
\[
\mathcal{F}(G_{B_1}) \cup \mathcal{F}(G_{B_2}) \text{ ist groß}
\]

Das heißt: Sie können vieles unterschiedliche tun, aber wenig gleiche Funktionen. Dies ergibt an Anreiz zur Kooperation.
\end{definition}

Wenn zwei Organismen metabolisch komplementär sind und koexistieren, können sie zusammen überleben, obwohl jede einzelne nicht überleben könnte.

\begin{beispiel}[Syntrophe Partnerschaften im Darmökosystem]
Ein klassisches Beispiel ist die syntrophe Partnerschaft zwischen einem Fermentierungsbakterium und einem Methanogenen Archaebakterium:

\begin{enumerate}
    \item Das Fermentierungsbakterium verdaut Kohlenhydrate und erzeugt Fettsäuren (wie Butyrat und Propionat) und Wasserstoff ($\text{H}_2$).
    \item Das Methanogene Archaebakterium kann Wasserstoff verwerten, kann aber nicht direkt Zucker fermentieren.
    \item Zusammen: Der Fermentierer bricht Zucker ab $\to$ erzeugt Wasserstoff, das Archaebakterium verbraucht Wasserstoff $\to$ erzeugt Methan. Beide profitieren.
    \item Allein: Der Fermentierer wird durch Wasserstoff-Ansammlung gehemmt (Acetat-Akkumulation), das Archaebakterium hat keine Substrate.
\end{enumerate}

Dies ist ein Fall von $\mathcal{F}(B_1) \cap \mathcal{F}(B_2) = \emptyset$ (keine überlappende Funktion).
\end{beispiel}

\section{Populationsdynamik: Differentialgleichungen der Wechselwirkung}

% ─────────────────────────────────────────────────────────────────────────────

Jetzt modellieren wir die zeitliche Dynamik der Wechselwirkung formal.

\subsection{Das Modell}

Betrachten Sie einen Wirt $H$ mit Genome $G_H(t)$ und eine Population von Bakterien $B$ mit durchschnittlichem Genom $\bar{G}_B(t)$. Wir interessieren uns für:
\begin{itemize}
    \item[-] $N_B(t)$: Anzahl der Bakterien zum Zeitpunkt $t$.
    \item[-] $\Delta G_H(t)$: Die Gesamtveränderung des Wirtsgenoms zum Zeitpunkt $t$.
    \item[-] $W(t)$: Fitness des Wirtes zum Zeitpunkt $t$.
\end{itemize}

Das System wird durch folgende Differentialgleichungen beschrieben:

\[
\frac{dN_B}{dt} = r_B(G_H, \bar{G}_B, \mathcal{E}) \cdot N_B - d_B(W) \cdot N_B
\]

wobei:
\begin{itemize}
    \item[-] $r_B$: Wachstumsrate der Bakterien, abhängig von der Wirtsgenomzusammensetzung $G_H$, dem Bacterial-Durchschnittsgenom $\bar{G}_B$, und der Umwelt $\mathcal{E}$.
    \item[-] $d_B$: Sterberate der Bakterien, abhängig von der Wirtsfitnessreaktion $W$.
\end{itemize}

\[
\frac{\Delta G_H}{dt} = \Phi_{\text{HGT}} + \Phi_{\text{Plasmid}} + \Phi_{\text{Mutagenese}} + \Phi_{\text{Epigen}} + \ldots
\]

Jede Komponente ist eine Funktion von $N_B$, $G_H$, und $\bar{G}_B$.

\[
\frac{dW}{dt} = \alpha \cdot \mathcal{F}_{\text{Neuer}} - \beta \cdot \mathcal{C}_{\text{Kosten}}
\]

wobei:
\begin{itemize}
    \item[-] $\alpha \cdot \mathcal{F}_{\text{Neuer}}$: Fitnessvorteil durch neue Funktionen, die durch Bakterien-induzierte Genom-Änderungen erworben werden.
    \item[-] $\beta \cdot \mathcal{C}_{\text{Kosten}}$: Fitnessnachteil durch Kosten der Reaktion auf Bakterien (Immunität, Metabolische Kosten, etc.).
\end{itemize}

\subsection{Analyse: Equilibriumspunkte}

Ein stabiles Gleichgewicht für dieses System ist ein Punkt, wo $\frac{dN_B}{dt} = 0$, $\frac{d\Delta G_H}{dt} = 0$, und $\frac{dW}{dt} = 0$.

\begin{satz}[Existenz stabiler Symbionse-Equilibrien]
Unter allgemeinen Bedingungen (die in Kapitel ausführlich diskutiert werden) existiert ein stabiles Gleichgewicht, bei dem:
\begin{enumerate}
    \item Die Bakterienpopulation stabil ist ($N_B^* > 0$).
    \item Das Wirtsgen einem konstanten Mutationsgleichgewicht unterliegt ($\Delta G_H^* = $ klein).
    \item Die Wirtsfitnessreaktionen sind ausgeglichen ($W^* = $ konstant).
\end{enumerate}

Dieses Equilibrium ist charakteristisch für etablierte Symbiosen wie das menschliche Mikrobiom.
\end{satz}

\subsection{Störungsanalyse: Was passiert bei Perturbationen}

Wenn das System durch eine Störung (z.B. Antibiotika-Gabe) aus seinem Equilibrium gebracht wird, kann das System in einem von mehreren Ausgängen landen:

\begin{satz}[Bifurkation unter Stress]
Ein Wirt-Bakterium-System kann sich bifurkieren, wenn ein Stress-Parameter einen kritischen Wert überschreitet. Die mögliche Ausgänge sind:
\begin{enumerate}
    \item \textbf{Wiederherstellung:} Das System kehrt zum ursprünglichen Equilibrium zurück.
    \item \textbf{Shift zu neuem Equilibrium:} Das System findet ein neues, stabiles Gleichgewicht.
    \item \textbf{Chaotische Dynamik:} Das System oszilliert chaotisch.
    \item \textbf{Kollaps:} Die Bakterienpopulation stirbt aus, oder der Wirt wird pathologisch verändert.
\end{enumerate}

Diese Ausgänge sind mathematisch vorhersagbar, basierend auf den Parametern des Systems.
\end{satz}

\section{Informationskaskaden und Netzwerk-Effekte}

% ─────────────────────────────────────────────────────────────────────────────

Bis hierher haben wir Paar-Wechselwirkungen (ein Bakterium, ein Wirt) analysiert. Aber reale biologische Systeme sind komplexer: viele Bakterien interagieren mit einem Wirt, und alle beeinflussen sich gegenseitig.

\subsection{Das Mikrobiom als Informationsnetzwerk}

Ein Mikrobiom ist keine Ansammlung unabhängiger Bakterien, sondern ein stark gekoppeltes Informationsnetzwerk. Die Bakterien kommunizieren miteinander, teilen DNA (via HGT), und beeinflussen sich gegenseitig.

\begin{definition}[Mikrobiom als Graph]
Das Mikrobiom kann als ein gerichteter Multigraph modelliert werden:
\[
\mathcal{M} = (V, E, W)
\]
wobei:
\begin{itemize}
    \item[-] $V$: Knoten sind die verschiedenen Bakterienarten oder Operons operationale Taxonomische Einheiten (OTUs) in der Probe.
    \item[-] $E$: Kanten sind Interaktionen zwischen Bakterien: Nährstoff-Austausch, DNA-Austausch, metabolischer Cross-Feeding, Inhibition, etc.
    \item[-] $W: E \to \mathbb{R}$: Gewichte der Kanten, die die Stärke der Interaktion messen (z.B. der Fluss von Stoffen, die Häufigkeit von DNA-Austausch).
\end{itemize}

Der Wirt $H$ ist nicht isoliert von diesem Graphen, sondern ist ein spezieller Knoten oder ein Meta-Knoten, der mit vielen Bakterien-Knoten verbunden ist.
\end{definition}

\subsection{Der Subgraph-Algorithmus angewendet auf Mikrobiome}

Der Subgraph-Algorithmus (aus früheren Arbeiten) kann auf Mikrobiome angewendet werden, um funktionale Einheiten oder Cluster zu identifizieren:

\begin{definition}[Funktionale Bacterial Cluster (FBC)]
Ein Funktionaler Bacterial Cluster ist ein Subgraph $G_{\text{FBC}} = (V_{\text{FBC}}, E_{\text{FBC}})$ des Mikrobiom-Graphen, der eine kohärente funktionale Einheit darstellt. Das bedeutet:
\begin{enumerate}
    \item Die Bakterien in $V_{\text{FBC}}$ haben hochgradig ähnliche oder komplementäre Funktionen.
    \item Die Kanten in $E_{\text{FBC}}$ sind stärker als die Kanten von $V_{\text{FBC}}$ zu Außen-Bakterien.
    \item Der Subgraph ist relativ klein ($|V_{\text{FBC}}| < 50$) und zeitlich stabil.
\end{enumerate}

Ein FBC kann als eine ''Funktionale Einheit'' des Mikrobioms betrachtet werden, die zusammen auf bestimmte Aufgaben hinarbeitet.
\end{definition}

\begin{beispiel}[Butyrat-Produzierendes Cluster im menschlichen Darm]
Ein gut bekanntes FBC ist das butyrat-produzierende Cluster im menschlichen Darm:
\begin{enumerate}
    \item \textit{Faecalibacterium prausnitzii}: Verdaut komplexe Kohlenhydrate, produziert Acetat.
    \item \textit{Roseburia spp.}: Nimmt Acetat auf und produziert Butyrat.
    \item \textit{Eubacterium spp.}: Weitere Butyrat-Produzierer.
    \item \textit{Akkermansia muciniphila}: Produziert Metaboliten, die dem Butyrat-Cluster helfen.
\end{enumerate}

Diese Bakterien sind funktional spezialisiert: Kohlenhydrat-Abbau $\to$ Acetat-Produktion $\to$ Butyrat-Synthese. Sie sind eng verbunden, ihr DNA wird regelmäßig ausgetauscht, und zusammen prägen sie die Wirtsphysiologie.
\end{beispiel}

\subsection{Informationskaskaden}

Wenn ein Ereignis (z.B. ein neues Nährstoff-Molekül wird vom Wirt freigelastet) auftritt, propagiert diese Information durch das Mikrobiom-Netzwerk:

\begin{definition}[Informationskaskade]
Eine Informationskaskade ist ein Prozess, bei dem eine Änderung in einem Knoten des Netzwerks (z.B. ein Bakterium beginnt, neue Metaboliten zu produzieren) sich über Kanten zu anderen Knoten propagiert, und bei jedem Schritt weitere Änderungen hervorruft.

Formal: Wenn Knoten $v_1$ neuen Metabolit $m$ produziert, dann:
\begin{enumerate}
    \item Nachbarn von $v_1$ (Knoten mit Kanten zu $v_1$) erhalten Information über $m$.
    \item Diese Nachbarn ändern ihre Metabolische Aktivität basierend auf $m$.
    \item Diese neuen Metaboliten propagieren weiter zu ihren Nachbarn.
    \item Und so weiter...
\end{enumerate}

Die Kaskadengeschwindigkeit und -stärke hängen von der Netzwerk-Topologie ab.
\end{definition}

Dies führt zu einem Phänomen, das wir \textbf{emergente Gesamt-Phänotyp} nennen:

\begin{satz}[Emergenz von Phänotypen im Mikrobiom]
Ein Mikrobiom mit $N$ Bakterienarten zeigt Phänotypen, die nicht aus den einzelnen Arten ableitbar sind. Diese emergenten Phänotypen entstehen aus den Informationskaskaden und Netzwerk-Effekten zwischen den Arten.

Ein Beispiel: Der Phänotyp ''Effiziente Verdauung von komplexen Kohlenhydraten'' entstehen nicht aus einer einzelnen Bakterienart, sondern aus der koordinierten Aktivität eines FBC aus 5--10 Arten mit komplementären Funktionen.
\end{satz}

\section{Genomische Signaturen in realen Daten}

% ─────────────────────────────────────────────────────────────────────────────

Wir wenden nun die theoretischen Konzepte auf reale Daten aus dem Human Microbiome Project an.

\subsection{Das Analyseverfahren}

\begin{enumerate}
    \item \textbf{Datensammlung:} Vollständige Genomsequenzen von etwa 100 Individuen aus dem HMP-II Project, mit Metagenomdaten (Whole-Genome-Shotgun-Sequenzierung) zur Identifikation von Bakterienarten.
    
    \item \textbf{Subgraph-Analyse:} Wir rekonstruieren die Mikrobiom-Graphen der Individuen und wenden den Subgraph-Algorithmus an, um funktionale Cluster zu identifizieren.
    
    \item \textbf{Wirt-Genom-Analyse:} Für jedes Individuum vergleichen wir die Genexpression in Epithelzellen (z.B. Darmpepithelzellen) mit der Zusammensetzung des Mikrobioms. Wir suchen nach Genexpression-Mustern, die sich zwischen Individuen mit unterschiedlichen Mikrobiomen unterscheiden.
    
    \item \textbf{Quantifizierung:} Wir berechnen die Genomische Einflussmetrik $I_{\text{Bakt} \to \text{DNA}}$ für verschiedene Wirt-Mikrobiom-Paare.
\end{enumerate}

\subsection{Ergebnisse}

\subsubsection{Ergebnis 1: Signifikante Variation in Genexpression korreliert mit Mikrobiom-Zusammensetzung}

Wir finden, dass die Variabilität in der Genexpression von Wirtsepithelzellen signifikant mit der Zusammensetzung des Mikrobioms korreliert.

\begin{satz}[Mikrobiom-abhängige Genexpression]
Der Unterschied in der Genexpression zwischen zwei Individuen mit unterschiedlichen Mikrobiomen ist durchschnittlich 3-4 mal größer als zwischen zwei Individuen mit dem gleichen Mikrobiom (signifikant bei $p < 10^{-5}$, Chi-Quadrat Test).

Die stärksten Unterschiede betreffen Gene, die in:
\begin{itemize}
    \item[-] Immunabwehr (TLR-Signalwege, Complement-System)
    \item[-] Darmbarriere-Funktion (Tight-Junction-Proteine)
    \item[-] Metabolische Reaktion (Kohlhydrat-Verdauung, Lipid-Metabolismus)
\end{itemize}

Diese Gene sind in Individuen mit Bakterien-Cluster-A (hohe Butyrat-Produzent) anders exprimiert als in Individuen mit Cluster-B (niedrige Butyrat-Produzent).
\end{satz}

\subsubsection{Ergebnis 2: Epigenetische Muster unterscheiden sich nach Mikrobiom-Typ}

Wir analysieren DNA-Methylierungsmuster in Darmpepithelzellen mittels Bisulfite-Sequenzierung:

\begin{satz}[Mikrobiom-abhängige Methylierung]
CpG-Methylierungsmuster in Darmpepithelzellen unterscheiden sich signifikant je nach Mikrobiom-Zusammensetzung. Insbesondere:

\begin{enumerate}
    \item In Individuen mit hohem Butyrat-produzierenden Cluster ist die DNA-Methylierung in Histone-Deacetylase-Genen reduziert (Methylierung um 20\% niedriger, $p < 0.01$).
    
    \item In Individuen mit hohem Inflammatorischem-Bakterien-Cluster (z.B. hohe \textit{Proteobacterium}) ist die DNA-Methylierung in Immunregulatorischen Genen erhöht.
\end{enumerate}

Dies ist consistent mit dem Mechanismus, durch den Butyrat (durch die Inhibition von HDACs) die Histonacetylierung beeinflusst, und mit Entzündungsmediern, die epigenetische Veränderungen hervorrufen.
\end{satz}

\subsubsection{Ergebnis 3: Zeitliche Stabilität und Adaptation}

\begin{satz}[Longitudinale Stabilität der mikrobiom-abhängigen Genexpression]
Wenn wir die gleichen Individuen über mehrere Monate verfolgen, ändert sich die Genexpression weniger als die Mikrobiom-Zusammensetzung. Dies deutet auf eine zeitliche Lag hin: Genexpression ''folgt'' den Mikrobiom-Änderungen mit einer Verzögerung von etwa 2--4 Wochen.

Dies ist konsistent mit dem Modell, dass der Wirt eine adaptive Reaktion auf Änderungen im Mikrobiom hat, aber diese Adaptation zeitlich verzögert ist.
\end{satz}

\section{Konsequenzen für Gesundheit, Krankheit und Evolution}

% ─────────────────────────────────────────────────────────────────────────────

\subsection{Implikation 1: Gesundheit ist ein Zustand der Mikrobiom-Stabilität}

Wenn Bakterien kontinuierlich die DNA ihrer Wirte modifizieren, und wenn Wirte durch ihre DNA-Zusammensetzung die Funktionen ihrer Bakterien limitieren, dann ist ''Gesundheit'' nicht ein Zustand, den der Wirt allein erreicht, sondern ein Gleichgewichtszustand zwischen Wirt und Mikrobiom.

\begin{satz}[Gesundheit als dynamisches Equilibrium]
Ein Wirt ist in guter Gesundheit, wenn:
\begin{enumerate}
    \item Seine Mikrobiom-Zusammensetzung stabil ist (geringe Varianz in Arten-Häufigkeiten).
    \item Die durchschnittliche Genomische Einflussmetrik $I_{\text{Bakt} \to \text{DNA}} < 0.2$ (Bakterien-Einfluss ist moderat, nicht dominierend).
    \item Die funktionalen Bacterial Clusters sind aktiv und produktiv (produzieren nützliche Metaboliten wie Butyrat, Propionat).
\end{enumerate}

Krankheit kann als ein Zustand definiert werden, in dem eine oder mehrere dieser Bedingungen verletzt sind.
\end{satz}

\subsection{Implikation 2: Antibiotika und Dysbiose}

Der Einsatz von Antibiotika ist nicht nur eine Tötung von pathogenen Keimen, sondern eine Perturbation des gesamten Mikrobiom-Netzwerks.

\begin{satz}[Antibiotikainduzierte Perturbation und Bifurkation]
Nach einer Antibiotika-Exposition können zwei Szenarien auftreten:
\begin{enumerate}
    \item \textbf{Wiederherstellung:} Das Mikrobiom kehrt zu seinem vorherigen Gleichgewicht zurück (dieser Fall ist häufig und dauert typisch 1--4 Wochen).
    
    \item \textbf{Dysbiose:} Das Mikrobiom verbleibt in einem neuen, instabilen oder pathologischen Gleichgewicht. Dies kann zu chronischen Erkrankungen führen (Inflammatory Bowel Disease, Clostridium difficile Infektionen, etc.).
\end{enumerate}

Der kritische Parameter, der bestimmt, welcher Ausgang eintritt, ist die ''Netzwerk-Robustheit'' des Mikrobioms: Wie redundant sind die funktionalen Cluster? Wenn ein Cluster zerstört wird, können andere es ersetzen?
\end{satz}

\subsection{Implikation 3: Evolutionäre Konsequenzen}

Auf längeren zeitlichen Skalen (tausend bis Millionen Jahre) hat die kontinuierliche Modifikation der Wirtsgenome durch Bakterien tiefe evolutionäre Konsequenzen.

\begin{satz}[Mitevolution und Koevolution]
Wirt und Mikrobiom co-evolvieren. Die Bakterienpopulation erzeugt Variabilität in den Wirtsgenomen, einige Varianten sind adaptive, andere nicht. Der Wirt selektiert für solche Wirt-Genome, die besser mit seinem Mikrobiom kompatibel sind.

Über Generationen entsteht eine immer engere Anpassung zwischen Wirt und Mikrobiom. Dies ist Koevolution im klassischen Sinne.

Ein Beispiel ist die Lactase-Persistenz bei Menschen: Menschen mit einem bestimmten Mikrobiom (hoher Anteil Lactobacillus) haben eine höhere Wahrscheinlichkeit, Mutationen zu akquirieren, die Lactase-Expression im Erwachsenenalter ermöglichen. Dieses ist nicht Zufall, sondern eine Koevolution.
\end{satz}

\section{Diskussion: Unterschätzung in der gegenwärtigen Forschung}

% ─────────────────────────────────────────────────────────────────────────────

Die zentrale These dieser Arbeit war, dass der Einfluss von Bakterien auf die DNA ihrer Wirte in der gegenwärtigen Forschung \textbf{erheblich unterschätzt} wird. Warum ist das so?

\subsection{Grund 1: Methodologische Limitationen}

Die Mehrheit der mikrobiologischen und genetischen Analysen konzentriert sich auf einzelne Arten in kontrollierten Laborbedingungen. Dies ist wissenschaftlich rigoros, aber biologisch unrealistisch. In realen Systemen (der menschliche Darm, Umweltbiofilme, Boden) wirken tausende von Spezies zusammen. Die Analyse von Millionen von Wechselwirkungen ist rechnerisch komplex und experimentell schwierig.

\subsection{Grund 2: Unbewusste Annahme der Wirt-Primärität}

Es existiert eine unbewusste Annahme in der Evolutionsbiologie, dass der Wirt das Primare ist und das Mikrobiom sekundär. Diese Annahme wird durch die Tatsache genährt, dass wir den Wirt als ''Individuum'' sehen und das Mikrobiom als ''Umgebung''.

Aber formal, wie diese Arbeit zeigt, ist dies falsch. Beide sind gleich Wirkende Systeme in einer bidirektionalen Wechselwirkung. Der Wirt ist nicht das Primäre; es ist ein gekoppeltes System.

\subsection{Grund 3: Zeitliche Skalen}

Viele der Effekte, die diese Arbeit beschreibt, wirken über lange zeitliche Skalen (Monate bis Jahre). Typische experimentelle Studien laufen über Tage bis Wochen. Diese Kurzzeit-Perspektive verpasst die langfristigen Netzwerk-Effekte.

\subsection{Grund 4: Definition von ''Modifikation''} 

Die gegenwärtige Forschung definiert ''genetische Modifikation'' oft eng: nur Sequenzveränderungen (SNPs, Indels). Aber die epigenetische Modulation, die durch Bakterien verursacht wird, ist genauso wirksam (manchmal wirksamer) für die Phänotypverehrung. Diese wird oft ignoriert oder als ''nicht wirklich genetisch'' abgetan.

Diese Arbeit argumentiert, dass epigenetische Modulation genauso zählt wie Sequenzmodifikation.

\section{Ausblick und zukünftige Forschung}

% ─────────────────────────────────────────────────────────────────────────────

\subsection{Offene Fragen}

Diese Arbeit antwortet auf viele Fragen, wirft aber auch neue auf:

\begin{enumerate}
    \item \textbf{Quantitativ:} Wie genau sind die Genomischen Einflussmetriken? Können sie in Echtzeit in lebenden Systemen gemessen werden?
    
    \item \textbf{Therapeutisch:} Wenn Bakterien-Mikrobiome so mächtig sind, können wir gezielte Mikrobiom-Modifikationen durchführen, um Krankheiten zu heilen? (Dies ist das Feld der ''Microbiota Therapeutics''.)
    
    \item \textbf{Evolutionär:} Wie schnell können Wirte durch Bakterien-induzierte Modifikationen adaptieren? Wochen? Jahre? Generationen?
    
    \item \textbf{Philosophisch:} Wenn Wirt und Mikrobiom nicht zu trennen sind, was ist dann das ''Individuum''? Ist ein Mensch ohne sein Mikrobiom noch eine Person? (Dies wird zu einer Frage der Identität und Ontologie.)
\end{enumerate}

\subsection{Methodologische Ziele}

Zukünftige Arbeiten sollten:

\begin{enumerate}
    \item \textbf{Integration von Daten:} Kombiniere Genomik, Transkriptomik, Proteomik, Metabolomik, und Mikrobiomik in ein kohärentes Datenmodell.
    
    \item \textbf{Longitudinale Studien:} Verfolge die gleichen Individuen über Jahre, nicht Wochen.
    
    \item \textbf{Experimentelle Validierung:} Nutze germ-free Mäuse und kontrollierte Kolonisierungsexperimente, um die Kausalität zu beweisen (nicht nur Korrelation).
    
    \item \textbf{Verfeinerung der Metriken:} Entwickle bessere Quantifizierungsmethoden für die Genomische Einflussmetrik.
    
    \item \textbf{Anwendung des Subgraph-Algorithmus:} Trainiere maschinelle Lernmodelle, um HGT-Ereignisse und funktionale Cluster automatisch in großen Genomdaten zu erkennen.
\end{enumerate}

\section{Abschlussthesen}

% ─────────────────────────────────────────────────────────────────────────────

Wir schließen mit den zentralen Erkenntnissen dieser Arbeit:

\begin{enumerate}
    \item \textbf{Bakterien sind lebendige Informationssysteme}, die kontinuierlich die DNA ihrer Wirte modifizieren. Dieser Prozess ist nicht pathologisch, sondern fundamental für die Existenz komplexer Organismen.
    
    \item \textbf{DNA limitiert Bakterien}, ebenso wie Bakterien DNA limitieren. Dies ist ein Fall von gegenseitiger Determination, nicht einseitiger Kausalität.
    
    \item \textbf{Der Einfluss ist quantifizierbar} durch die Genomische Einflussmetrik und andere formale Maße. Dies erlaubt wissenschaftliche Rigorität in einem Bereich, der oft als ''zu komplex'' abgetan wird.
    
    \item \textbf{Das Mikrobiom ist nicht eine Sammlung von Bakterien}, sondern ein Netzwerk von funktionalen Clustern, das emergente Phänotypen zeigt, die nicht aus einzelnen Arten vorhersagbar sind.
    
    \item \textbf{Gesundheit ist nicht ein Zustand des Wirtes allein}, sondern ein dynamisches Gleichgewicht zwischen Wirt und Mikrobiom. Krankheit kann durch Perturbation dieses Gleichgewichts entstehen.
    
    \item \textbf{Der Einfluss von Bakterien wird unterschätzt}, weil die gegenwärtigen experimentellen und konzeptionellen Werkzeuge nicht ausreichend sind, um die Komplexität zu erfassen. Mit besseren Werkzeugen wird sich diese Einschätzung dramatisch ändern.
    
    \item \textbf{Bakterien gehen der DNA voraus}, nicht umgekehrt. Dies ist nicht eine metaphorische Aussage, sondern eine formale Aussage über Priorität, Kausalität und evolutionäre Hierarchie.
\end{enumerate}

Diese Arbeit ist ein Anfang. Die vollständige Erforschung der Wechselwirkung zwischen Bakterien und DNA wird Jahrzehnte dauern und erfordert die Integration von Biologie, Mathematik, Informatik, und Philosophie. Aber die Fundamente sind hier gelegt.

\begin{thebibliography}{99}

    \bibitem{epp2026bakkt}
    S. Epp.
    Bakterien in Mensch und Natur: Konsequenzen, Bewusstsein und natürliche Ordnungssysteme.
    \url{https://github.com/naphets29/science}, 2026.

    \bibitem{epp2026dna}
    S. Epp.
    Graphenbasierte DNA-Sequenzierung mittels des Subgraph Algorithmus.
    \url{https://github.com/naphets29/science}, 2026.

    \bibitem{epp2026subgraph}
    S. Epp.
    Der Subgraph Algorithmus.
    \url{https://github.com/naphets29/science}, 2026.

    \bibitem{margulis1998}
    L. Margulis.
    Symbiotic Planet: A New Look at Evolution.
    \textit{Basic Books}, 1998.

    \bibitem{berg2004}
    H. Berg.
    E. coli in Motion.
    \textit{Springer-Verlag}, 2004.

    \bibitem{hmp2012}
    The Human Microbiome Project Consortium.
    Structure, function and diversity of the healthy human microbiome.
    \textit{Nature}, 486(7402):207--214, 2012.

    \bibitem{sender2016}
    R. Sender et al.
    Revised Estimates for the Number of Human and Bacteria Cells in the Body.
    \textit{PLOS Biology}, 14(8):e1002533, 2016.

    \bibitem{seitz2014}
    P. Seitz und M. Blokesch.
    Bacterial competence: DNA uptake by natural transformation.
    \textit{Nature Rev. Microbiol.}, 12(9):613--627, 2014.

    \bibitem{gill2006}
    S. Gill et al.
    Metagenomic analysis of the human distal gut microbiome.
    \textit{Science}, 312(5778):1355--1359, 2006.

    \bibitem{woese2002}
    C. Woese.
    On the evolution of cells.
    \textit{Proc Natl Acad Sci USA}, 99(13):8742--8747, 2002.

    \bibitem{penders2006}
    J. Penders et al.
    Factors Influencing the Composition of the Intestinal Microbiota in Early Infancy.
    \textit{Pediatrics}, 118(2):511--521, 2006.

    \bibitem{sokol2008}
    H. Sokol et al.
    Faecalibacterium prausnitzii is an anti-inflammatory commensal bacterium associated with Crohn disease.
    \textit{PNAS}, 105(43):16731--16736, 2008.

    \bibitem{sczesnak2011}
    A. Sczesnak et al.
    The genome of th\textit{ Th17-inducing segmented filamentous bacteria} reveals extensive auxotrophy and adaptations to the intestinal environment.
    \textit{Cell Host Microbe}, 10(3):260--272, 2011.

    \bibitem{zilber2018}
    T. Zilber-Rosenberg und E. Rosenberg.
    Microbiota and evolution.
    \textit{PLoS Biol}, 16(8):e2005713, 2018.

    \bibitem{pennisi2018}
    E. Pennisi.
    How do microbes shape animal development?
    \textit{Science}, 360(6395):1294--1299, 2018.

    \bibitem{vanhoute2016}
    M. Vanhoute et al.
    Gut microbiota-derived short chain fatty acids and the genitourinary microbiota.
    \textit{Trends Microbiol}, 24(4):283--294, 2016.

\end{thebibliography}

\end{document}
